\begin{abstract}
% Five sentences: (1) what this study shows, (2) why the question matters, (3) how it is tested,
% (4) the evidence, (5) the most important number. Every claim matches a test-set result in
% Section~\ref{sec:results}.
Diffusion-generated images help few-shot image classifiers under narrower conditions than their low cost
suggests: on CIFAR-100 they help pretrained classifiers and add to conventional augmentation, but they lose to it
when they replace it from random initialization, and on a fine-grained and a medical dataset they do not beat
it. Generated images are cheap and come with labels, but they differ from real images
in ways that standard quality measures miss, and earlier studies compared them under mixed protocols. We train
ResNet-50, ViT-B/16 and ConvNeXt-T from ImageNet weights and from random initialization on CIFAR-100 with 5 to 50
images per class, compare four generated sets with six conventional augmentations at matched training cost,
paired over three to five seeds, and repeat the study on Flowers-102 and HAM10000. With pretrained classifiers,
Stable Diffusion class-prompt images raise test accuracy by up to 14.2 points at 5 images per class, while from
random initialization they lose to RandAugment in 9 of 10 cells; FID and recall rank the generated sets in the
opposite order to their usefulness, and on Flowers-102 and HAM10000 neither class-prompt nor LoRA images, with RandAugment,
beat the best conventional augmentation in any cell. The largest gain, 14.9 points over the best conventional augmentation chosen on
validation (ResNet-50, 5 images per class), comes from adding five generated images per real image on top of
RandAugment, so generated images are best used to extend conventional augmentation, after a validation pilot on
the target data.
\end{abstract}
